\subsection{Biology-Based Algorithms}

% =========================================================
\subsubsection{Particle Swarm Optimization (PSO)}
% =========================================================

\textbf{Introduction}

Particle Swarm Optimization (PSO) is a population-based metaheuristic proposed by Kennedy and
Eberhart in 1995, inspired by the collective movement of bird flocks and fish schools. Each
candidate solution is modeled as a \textit{particle} that navigates the search space by
combining its own memory (personal best position) with social knowledge shared across the
entire swarm (global best position). The implementation supports both continuous optimization
and discrete combinatorial problems such as the Travelling Salesman Problem (TSP).

Default hyperparameters used in this implementation:
\begin{center}
\begin{tabular}{lll}
\hline
\textbf{Parameter} & \textbf{Symbol} & \textbf{Default value} \\
\hline
Population size          & $N$          & 30   \\
Number of iterations     & $G$          & 100  \\
Max inertia weight       & $w_{\max}$   & 0.9  \\
Min inertia weight       & $w_{\min}$   & 0.4  \\
Cognitive coefficient    & $c_1$        & 2.05 \\
Social coefficient       & $c_2$        & 2.05 \\
\hline
\end{tabular}
\end{center}

\textbf{Operating Principle}

Each particle $i$ at iteration $t$ carries three pieces of state:
\begin{itemize}
  \item \textbf{Position} $\mathbf{x}_i^{(t)}$: the current candidate solution in the
        $D$-dimensional search space.
  \item \textbf{Velocity} $\mathbf{v}_i^{(t)}$: the direction and step size for the next
        positional update.
  \item \textbf{Personal best} $\mathbf{p}_i^{\text{best}}$: the best position particle $i$
        has ever visited, together with its cost $f(\mathbf{p}_i^{\text{best}})$.
\end{itemize}
A single \textbf{global best} $\mathbf{g}^{\text{best}}$ records the lowest-cost position
found by any particle across all iterations. At every step, each particle is pulled toward
both its own personal best and the global best, while inertia carries it along its previous
trajectory. The inertia weight $w$ decays linearly iteration by iteration, shifting the swarm
from broad exploration early on to fine-grained exploitation by the final iteration.

For TSP, positions are represented as continuous vectors in $[0,1]^n$ and converted to valid
permutations via \texttt{pos\_to\_tour} (index-sorting). Velocities are clamped to
$[-0.4,\,0.4]$ per dimension to maintain permutation stability.

\textbf{Update Mechanism}

At each iteration $t$, the inertia weight is first computed:
\[
w^{(t)} = w_{\max} - \frac{w_{\max} - w_{\min}}{G - 1}\cdot t
\]

Then, for every particle $i$, independent random vectors $\mathbf{r}_1,\mathbf{r}_2
\sim \mathcal{U}(0,1)^D$ are sampled and the velocity is updated:
\[
\mathbf{v}_i^{(t+1)} =
  \underbrace{w^{(t)}\,\mathbf{v}_i^{(t)}}_{\text{inertia}}
+ \underbrace{c_1\,\mathbf{r}_1 \odot \!\left(\mathbf{p}_i^{\text{best}} - \mathbf{x}_i^{(t)}\right)}_{\text{cognitive pull}}
+ \underbrace{c_2\,\mathbf{r}_2 \odot \!\left(\mathbf{g}^{\text{best}} - \mathbf{x}_i^{(t)}\right)}_{\text{social pull}}
\]

The position is then updated and clipped to the feasible domain:
\[
\mathbf{x}_i^{(t+1)} = \mathrm{clip}\!\left(\mathbf{x}_i^{(t)} + \mathbf{v}_i^{(t+1)},\;
                        \mathbf{lb},\; \mathbf{ub}\right)
\]

Personal and global bests are refreshed whenever a strictly lower cost is found.
The fitness value returned to the caller is transformed as:
\[
\text{fitness} =
\begin{cases}
\dfrac{1}{f(\mathbf{x}) + 1} & \text{if } f(\mathbf{x}) \ge 0 \quad (\text{continuous})\\[6pt]
f(\mathbf{x})                 & \text{(TSP, raw tour cost)}
\end{cases}
\]

\textbf{Algorithm}

\begin{algorithm}[H]
\caption{Particle Swarm Optimization}
\label{alg:pso}
\begin{algorithmic}[1]
\REQUIRE $f(\cdot)$, bounds $[\mathbf{lb},\mathbf{ub}]$, $N$, $G$,
         $w_{\max}$, $w_{\min}$, $c_1$, $c_2$
\ENSURE Best solution $\mathbf{g}^{\text{best}}$

\STATE $\mathbf{x}_i \sim \mathcal{U}(\mathbf{lb},\mathbf{ub})$,\;
       $\mathbf{v}_i \sim \mathcal{U}(-1,1)^D$ \quad for all $i = 1,\ldots,N$
\STATE $\mathbf{p}_i^{\text{best}} \leftarrow \mathbf{x}_i$,\;
       $\text{cost}_i \leftarrow f(\mathbf{x}_i)$ \quad for all $i$
\STATE $\mathbf{g}^{\text{best}} \leftarrow \mathbf{x}_{i^*}$ where
       $i^* = \arg\min_i \text{cost}_i$

\FOR{$t = 0$ \textbf{to} $G-1$}
  \STATE $w \leftarrow w_{\max} - (w_{\max}-w_{\min})\cdot t\,/\,(G-1)$
  \FOR{$i = 1$ \textbf{to} $N$}
    \STATE $\mathbf{r}_1, \mathbf{r}_2 \sim \mathcal{U}(0,1)^D$
    \STATE $\mathbf{v}_i \leftarrow w\,\mathbf{v}_i
           + c_1\,\mathbf{r}_1 \odot (\mathbf{p}_i^{\text{best}}-\mathbf{x}_i)
           + c_2\,\mathbf{r}_2 \odot (\mathbf{g}^{\text{best}}-\mathbf{x}_i)$
    \STATE $\mathbf{x}_i \leftarrow \mathrm{clip}(\mathbf{x}_i + \mathbf{v}_i,\,
           \mathbf{lb},\,\mathbf{ub})$
    \STATE $c \leftarrow f(\mathbf{x}_i)$
    \IF{$c < \text{cost}_i$}
      \STATE $\mathbf{p}_i^{\text{best}} \leftarrow \mathbf{x}_i$;\;
             $\text{cost}_i \leftarrow c$
      \IF{$c < f(\mathbf{g}^{\text{best}})$}
        \STATE $\mathbf{g}^{\text{best}} \leftarrow \mathbf{x}_i$
      \ENDIF
    \ENDIF
  \ENDFOR
\ENDFOR
\RETURN $\mathbf{g}^{\text{best}}$
\end{algorithmic}
\end{algorithm}

\textbf{Complexity Analysis}

\begin{itemize}
  \item \textbf{Time Complexity:} $O(G \cdot N \cdot D)$ per run (excluding the cost of $f$).
        Each of the $G$ iterations updates $N$ particles, each requiring $O(D)$ arithmetic
        operations for the velocity and position updates. Fitness evaluation adds $O(G\cdot N
        \cdot F)$ where $F$ is the cost of a single call to $f$.
  \item \textbf{Space Complexity:} $O(N \cdot D)$ to store positions $\mathbf{x}_i$,
        velocities $\mathbf{v}_i$, and personal bests $\mathbf{p}_i^{\text{best}}$ for all
        $N$ particles.
\end{itemize}

\textbf{Advantages}

\begin{itemize}
  \item \textbf{Simple and low-overhead:} Only six hyperparameters are required
        ($N$, $G$, $w_{\max}$, $w_{\min}$, $c_1$, $c_2$), and the per-iteration cost is
        linear in both $N$ and $D$.
  \item \textbf{Fast convergence on smooth landscapes:} The linearly decaying inertia weight
        shifts the swarm from global exploration ($w\approx 0.9$) to local exploitation
        ($w\approx 0.4$) automatically, making PSO highly effective on unimodal and
        near-convex functions such as Sphere and Ackley.
  \item \textbf{Derivative-free:} Only objective function evaluations are required; no
        gradient information is needed.
  \item \textbf{Unified continuous/discrete interface:} The \texttt{pos\_to\_tour}
        (index-sorting) encoding allows the same velocity-update rule to handle both
        continuous and TSP problems without structural changes to the algorithm.
\end{itemize}

\textbf{Disadvantages}

\begin{itemize}
  \item \textbf{Premature convergence on multimodal landscapes:} Because all particles share
        a single global best, the swarm tends to collapse around one attractor. On heavily
        multimodal functions like Rastrigin, PSO is frequently outperformed by methods that
        maintain greater population diversity (e.g., GA).
  \item \textbf{TSP velocity instability without clamping:} In the discrete encoding,
        unbounded velocities cause positions to leave $[0,1]^n$ and produce inconsistent
        permutations. The $[-0.4, 0.4]$ clamp resolves this but reduces the effective step
        size and slows exploration.
  \item \textbf{Stagnation near the end:} As $w$ approaches $w_{\min}=0.4$, the inertia term
        shrinks and the swarm loses momentum. If $\mathbf{g}^{\text{best}}$ has not been
        updated for many iterations, the remaining movement is dominated by small cognitive
        and social corrections and the search effectively halts.
  \item \textbf{Parameter sensitivity:} The choice of $c_1=c_2=2.05$ satisfies the
        constriction-factor stability condition but may not be optimal for every problem;
        re-tuning is often needed for best performance.
\end{itemize}

% =========================================================
\subsubsection{Ant Colony Optimization (ACO)}
% =========================================================

\textbf{Introduction}

Ant Colony Optimization (ACO) is a population-based metaheuristic proposed by Marco Dorigo in
1992, inspired by the pheromone-based foraging behavior of real ant colonies. Shorter paths
receive disproportionately more pheromone deposits per unit time, creating a self-reinforcing
feedback loop that gradually concentrates traffic on high-quality routes. This implementation
provides two distinct solvers selected automatically based on the problem type: a
\textit{discrete solver} for graph problems (TSP) that maintains a pheromone matrix and builds
probabilistic tours, and a \textit{continuous solver} (ACO$_\mathbb{R}$) that maintains a
ranked solution archive and samples new candidates from Gaussian kernels.

Default hyperparameters:
\begin{center}
\begin{tabular}{lll}
\hline
\textbf{Parameter} & \textbf{Symbol} & \textbf{Default value} \\
\hline
Ant / sample count      & $N$            & 10   \\
Number of iterations    & $G$            & 100  \\
Pheromone influence     & $\alpha$       & 1.0  \\
Heuristic influence     & $\beta$        & 2.0  \\
Evaporation rate        & $\rho$         & 0.5  \\
Pheromone constant      & $Q$            & 100  \\
Archive size            & $k$            & 50   \\
Convergence speed       & $\xi$          & 0.85 \\
\hline
\end{tabular}
\end{center}

\textbf{Operating Principle}

\textbf{Discrete mode (TSP).}
Each ant starts from a random city and constructs a complete tour by iteratively selecting the
next unvisited city through roulette-wheel selection. The selection probability for city $j$
from city $i$ is proportional to a combination of pheromone strength $\tau_{ij}$ and inverse
distance (heuristic) $\eta_{ij}=1/d_{ij}$. After all $N$ ants complete their tours, the
pheromone matrix is first decayed globally by factor $(1-\rho)$ and then reinforced on every
edge traversed, with deposit proportional to $Q/L_k$ where $L_k$ is the length of ant $k$'s
tour. Both directions of each edge receive the same deposit, reflecting the undirected nature
of TSP.

\textbf{Continuous mode (ACO$_\mathbb{R}$).}
The algorithm initialises an archive of $k$ solutions drawn uniformly from the search space.
Solutions are ranked by objective value and each rank $r \in \{0,\ldots,k-1\}$ is assigned a
Gaussian weight:
\[
w_r = \frac{1}{k\,\xi\,\sqrt{2\pi}}\exp\!\left(-\frac{r^2}{2\,\xi^2\,k^2}\right)
\]
Weights are normalised to sum to one. At each iteration $N$ new candidates are generated: for
each new solution, a source archive index is drawn according to the normalised weights, and
each dimension is sampled from a Gaussian centred at the corresponding archive value with a
standard deviation computed as a weighted average of absolute deviations across the archive
(floored at $1\%$ of the domain width). The archive is updated by merging old and new
solutions and retaining the $k$ best.

\textbf{Update Mechanism}

\textbf{Discrete ACO --- edge selection probability:}
\[
p_{ij} = \frac{[\tau_{ij}]^\alpha\,[\eta_{ij}]^\beta}
              {\displaystyle\sum_{l\in\mathcal{U}}[\tau_{il}]^\alpha\,[\eta_{il}]^\beta},
\quad j\in\mathcal{U}
\]
where $\mathcal{U}$ is the set of unvisited cities. Selection is performed by
\texttt{\_roulette\_select}, which normalises probabilities to unity before drawing from a
multinomial (using \texttt{numpy.random.choice}), eliminating floating-point bias.

\textbf{Discrete ACO --- pheromone update (applied after all ants finish):}
\[
\tau_{ij} \;\leftarrow\; (1-\rho)\,\tau_{ij}
         + \sum_{k:\,(i,j)\in\text{tour}_k} \frac{Q}{L_k}
\]
Both $\tau_{ij}$ and $\tau_{ji}$ receive the deposit because TSP edges are symmetric.

\textbf{Continuous ACO --- new solution generation for dimension $d$:}
\[
\sigma_d = \xi\sum_{i=0}^{k-1} w_i\,\bigl|s_i^{(d)} - s_{\text{src}}^{(d)}\bigr|,
\quad \sigma_d \leftarrow \max\!\left(\sigma_d,\; 0.01\,(ub_d - lb_d)\right)
\]
\[
x_d^{\text{new}} \sim \mathcal{N}\!\left(s_{\text{src}}^{(d)},\,\sigma_d^2\right),
\quad \text{clipped to } [lb_d, ub_d]
\]
where $s_{\text{src}}$ is the archive solution selected by the pheromone weights.

\textbf{Algorithm}

\begin{algorithm}[H]
\caption{ACO --- Discrete solver (TSP)}
\label{alg:aco-discrete}
\begin{algorithmic}[1]
\REQUIRE $f(\cdot)$, distance matrix $D$, $n$ cities, $N$, $G$, $\alpha$, $\beta$, $\rho$, $Q$
\ENSURE Best tour $\mathbf{x}^*$, best cost $f^*$

\STATE $\tau_{ij} \leftarrow 1$ for all $(i,j)$;\quad
       $\mathbf{x}^* \leftarrow \text{None}$;\quad $f^* \leftarrow \infty$
\FOR{$g = 1$ \textbf{to} $G$}
  \FOR{$k = 1$ \textbf{to} $N$} \COMMENT{Each ant builds one tour}
    \STATE $c \leftarrow \text{random city}$;\quad
           $\text{tour} \leftarrow [c]$;\quad
           $\mathcal{U} \leftarrow \{0,\ldots,n-1\}\setminus\{c\}$
    \WHILE{$\mathcal{U} \neq \emptyset$}
      \STATE $\eta_{cj} \leftarrow 1\,/\,\max(D_{cj},\,10^{-3})$ for all $j\in\mathcal{U}$
      \STATE $p_{cj} \leftarrow \tau_{cj}^\alpha\,\eta_{cj}^\beta\;/\;
             \sum_{l\in\mathcal{U}}\tau_{cl}^\alpha\,\eta_{cl}^\beta$
      \STATE $c \leftarrow \texttt{roulette\_select}(p)$;\quad
             $\text{tour} \leftarrow \text{tour}\,{+\!\!+}\,[c]$;\quad
             $\mathcal{U} \leftarrow \mathcal{U}\setminus\{c\}$
    \ENDWHILE
    \STATE $L_k \leftarrow f(\text{tour})$
    \IF{$L_k < f^*$} \STATE $f^*\leftarrow L_k$;\; $\mathbf{x}^*\leftarrow\text{tour}$
    \ENDIF
  \ENDFOR
  \STATE $\tau_{ij} \leftarrow (1-\rho)\,\tau_{ij}$ for all $(i,j)$ \COMMENT{Global evaporation}
  \FOR{each ant $k$}
    \STATE $\delta \leftarrow Q\,/\,L_k$
    \FOR{each edge $(i,j)$ in $\text{tour}_k$}
      \STATE $\tau_{ij} \mathrel{+}= \delta$;\quad $\tau_{ji} \mathrel{+}= \delta$
    \ENDFOR
  \ENDFOR
\ENDFOR
\RETURN $\mathbf{x}^*$, $f^*$
\end{algorithmic}
\end{algorithm}

\begin{algorithm}[H]
\caption{ACO --- Continuous solver (ACO$_\mathbb{R}$)}
\label{alg:aco-cont}
\begin{algorithmic}[1]
\REQUIRE $f(\cdot)$, bounds $[\mathbf{lb},\mathbf{ub}]$, $N$, $G$, $k$, $\xi$
\ENSURE Best solution $\mathbf{x}^*$

\STATE Initialise archive $\mathcal{A} \leftarrow$ $k$ random solutions; sort by $f$
\STATE $\mathbf{x}^* \leftarrow \mathcal{A}[0].\mathbf{s}$;\quad
       $f^* \leftarrow f(\mathbf{x}^*)$
\FOR{$g = 1$ \textbf{to} $G$}
  \STATE Compute $w_r \propto \exp\!\bigl(-r^2/(2\xi^2 k^2)\bigr)$; normalise \COMMENT{Rank weights}
  \FOR{sample $= 1$ \textbf{to} $N$}
    \STATE $\text{src} \leftarrow \texttt{weighted\_choice}(w)$
    \FOR{$d = 1$ \textbf{to} $D$}
      \STATE $\sigma_d \leftarrow \max\!\Bigl(\xi\!\sum_i w_i\,|s_i^{(d)}-s_{\text{src}}^{(d)}|,\;
             0.01(ub_d-lb_d)\Bigr)$
      \STATE $x_d \leftarrow \mathrm{clip}\bigl(\mathcal{N}(s_{\text{src}}^{(d)},\sigma_d^2),\,
             lb_d,\,ub_d\bigr)$
    \ENDFOR
    \IF{$f(\mathbf{x}) < f^*$}
      \STATE $f^* \leftarrow f(\mathbf{x})$;\; $\mathbf{x}^* \leftarrow \mathbf{x}$
    \ENDIF
  \ENDFOR
  \STATE $\mathcal{A} \leftarrow \text{top-}k\text{ of }\mathcal{A}\cup\{\text{new solutions}\}$
\ENDFOR
\RETURN $\mathbf{x}^*$
\end{algorithmic}
\end{algorithm}

\textbf{Complexity Analysis}

\textbf{Discrete ACO (TSP):}
\begin{itemize}
  \item \textbf{Time Complexity:} $O(G\cdot N\cdot n^2)$. Each of the $N$ ants visits $n$
        cities and at each step computes selection probabilities over $O(n)$ remaining
        candidates. The pheromone update is $O(N\cdot n)$ per iteration, dominated by tour
        construction.
  \item \textbf{Space Complexity:} $O(n^2)$ for the pheromone matrix $\boldsymbol{\tau}$ and
        the distance matrix.
\end{itemize}

\textbf{Continuous ACO (ACO$_\mathbb{R}$):}
\begin{itemize}
  \item \textbf{Time Complexity:} $O(G\cdot(N\cdot k\cdot D + k\log k))$. Generating each of
        the $N$ new solutions costs $O(k\cdot D)$ (computing $\sigma_d$ requires summing over
        all $k$ archive members for each of $D$ dimensions). Sorting the merged archive costs
        $O(k\log k)$.
  \item \textbf{Space Complexity:} $O(k\cdot D)$ for the solution archive.
\end{itemize}

\textbf{Advantages}

\begin{itemize}
  \item \textbf{Purpose-built for TSP:} The pheromone matrix lives directly on graph edges,
        so every deposit immediately encodes route quality; probabilistic tour construction
        guarantees feasible permutations without any repair step.
  \item \textbf{Normalised roulette-wheel selection:} The \texttt{\_roulette\_select}
        implementation explicitly normalises probabilities with NumPy before sampling,
        eliminating the floating-point rounding errors that accumulate in naive cumulative-sum
        approaches.
  \item \textbf{Adaptive bandwidth in continuous mode:} The per-dimension standard deviation
        $\sigma_d$ shrinks as the archive converges, automatically narrowing the search to
        exploit promising regions without requiring a hand-tuned cooling schedule.
  \item \textbf{Symmetric pheromone update:} Depositing on both $\tau_{ij}$ and $\tau_{ji}$
        after each tour correctly reflects the undirected structure of TSP, doubling
        reinforcement speed on good edges.
\end{itemize}

\textbf{Disadvantages}

\begin{itemize}
  \item \textbf{Quadratic tour-construction cost:} With $N$ ants and $n$ cities per tour, the
        discrete solver scales as $O(N\cdot n^2)$ per iteration, becoming prohibitive for
        large instances.
  \item \textbf{Stagnation through pheromone concentration:} If $\rho$ is set too small,
        pheromone on early good edges accumulates rapidly and locks all ants into a sub-optimal
        route before the colony has explored enough of the space.
  \item \textbf{Slow per-sample cost in continuous mode:} Computing $\sigma_d$ requires
        $O(k)$ operations \emph{per dimension per sample}, making ACO$_\mathbb{R}$ noticeably
        slower per iteration than PSO or Cuckoo Search at the same population size.
  \item \textbf{Many coupled parameters:} The discrete mode alone has four interacting
        hyperparameters ($\alpha$, $\beta$, $\rho$, $Q$); their joint effect on convergence
        speed versus solution quality is non-trivial to predict, requiring careful tuning or
        grid search.
\end{itemize}

% =========================================================
\subsubsection{Firefly Algorithm (FA)}
% =========================================================

\textbf{Introduction}

The Firefly Algorithm (FA) is a swarm-based metaheuristic proposed by Xin-She Yang in 2008,
inspired by the bioluminescent signalling of fireflies. Every firefly represents a candidate
solution; its \textit{brightness} is inversely proportional to the objective value, so a
brighter firefly corresponds to a better solution. The core movement rule is that a dimmer
firefly is attracted to --- and moves toward --- its single \textit{most attractive} brighter
neighbour (the one with the highest attractiveness $\beta$, not the globally brightest).
Fireflies with no brighter neighbour perform an independent random walk. This implementation
handles both continuous and TSP problems using the same movement equations; for TSP, positions
live in $[0,1]^n$ and are mapped to tours via index-sorting (\texttt{pos\_to\_tour}).

Default hyperparameters:
\begin{center}
\begin{tabular}{lll}
\hline
\textbf{Parameter} & \textbf{Symbol} & \textbf{Default value} \\
\hline
Population size        & $N$          & 50   \\
Number of iterations   & $G$          & 100  \\
Initial step size      & $\alpha_0$   & 0.5  \\
Base attractiveness    & $\beta_0$    & 1.0  \\
Light absorption coeff.& $\gamma$     & 1.0  \\
\hline
\end{tabular}
\end{center}

\textbf{Operating Principle}

The algorithm maintains a population of $N$ fireflies. At each iteration:
\begin{enumerate}
  \item The randomisation step size decays geometrically: $\alpha^{(t)} = \alpha_0 \cdot
        0.97^t$, gradually narrowing the perturbation radius.
  \item A copy of the current population is created (\texttt{new\_fireflies}). Each firefly
        $i$ scans all other fireflies $j$ that are brighter (lower cost). Among these it
        selects the \textit{single} most attractive neighbour --- the one with the highest
        $\beta_0 e^{-\gamma_{\text{eff}}\,r_{ij}^2}$ --- and moves toward it plus a random
        perturbation.
  \item If no brighter neighbour exists (firefly is the current best or locally dominant),
        it performs a larger random walk with amplitude $2\alpha$.
  \item After all moves, fitness is re-evaluated and the global best is updated.
\end{enumerate}
For continuous problems, $\gamma$ is normalised by the squared Euclidean diameter of the
search domain to make attractiveness scale-invariant. For TSP, the diameter is approximated as
$\sqrt{n}$ (positions in $[0,1]^n$).

\textbf{Update Mechanism}

\textbf{Scale-normalised absorption coefficient:}
\[
\gamma_{\text{eff}} =
\begin{cases}
\gamma\,/\,\|\mathbf{ub}-\mathbf{lb}\|_2^2 & \text{continuous}\\[4pt]
\gamma\,/\,n                               & \text{TSP, } d\approx\sqrt{n}
\end{cases}
\]

\textbf{Attractiveness of firefly $j$ as seen by firefly $i$:}
\[
\beta_{ij} = \beta_0\,\exp\!\bigl(-\gamma_{\text{eff}}\,r_{ij}^2\bigr),
\quad r_{ij} = \|\mathbf{x}_i - \mathbf{x}_j\|_2
\]

\textbf{Decaying step size:}
\[
\alpha^{(t)} = \alpha_0 \cdot 0.97^{t}
\]

\textbf{Position update} (firefly $i$ moves toward its best attractor $j^*$):
\[
j^* = \arg\max_{\{j\,:\,f(\mathbf{x}_j)<f(\mathbf{x}_i)\}}\beta_{ij}
\]
\[
\mathbf{x}_i^{\text{new}} =
\begin{cases}
\mathrm{clip}\!\left(\mathbf{x}_i + \beta_{ij^*}(\mathbf{x}_{j^*}-\mathbf{x}_i)
                     + \alpha^{(t)}\bigl(\mathbf{u}-0.5\bigr),\,\mathbf{lb},\,\mathbf{ub}\right)
  & \text{if } j^* \text{ exists}\\[6pt]
\mathrm{clip}\!\left(\mathbf{x}_i + 2\alpha^{(t)}\bigl(\mathbf{u}-0.5\bigr),\,
                     \mathbf{lb},\,\mathbf{ub}\right)
  & \text{otherwise (random walk)}
\end{cases}
\]
where $\mathbf{u}\sim\mathcal{U}(0,1)^D$.

The fitness transformation returned to the caller is:
\[
\text{fitness} =
\begin{cases}
1/(f(\mathbf{x})+1) & f(\mathbf{x})\ge 0 \text{ (continuous)}\\
f(\mathbf{x})       & \text{TSP (raw cost)}
\end{cases}
\]

\textbf{Algorithm}

\begin{algorithm}[H]
\caption{Firefly Algorithm}
\label{alg:fa}
\begin{algorithmic}[1]
\REQUIRE $f(\cdot)$, bounds, $N$, $G$, $\alpha_0$, $\beta_0$, $\gamma$
\ENSURE Best solution $\mathbf{x}^*$

\STATE $\mathbf{x}_i \sim \mathcal{U}(\mathbf{lb},\mathbf{ub})$ for $i=1,\ldots,N$
\STATE $\text{fit}_i \leftarrow f(\mathbf{x}_i)$;\quad
       $\mathbf{x}^* \leftarrow \mathbf{x}_{i^*}$, $f^*\leftarrow \text{fit}_{i^*}$
       where $i^*=\arg\min_i \text{fit}_i$
\STATE Compute $\gamma_{\text{eff}}$

\FOR{$t = 0$ \textbf{to} $G-1$}
  \STATE $\alpha \leftarrow \alpha_0\cdot 0.97^{t}$
  \STATE $\mathbf{x}^{\text{new}} \leftarrow \mathbf{x}$  \COMMENT{Copy current positions}
  \FOR{$i = 1$ \textbf{to} $N$}
    \STATE $j^* \leftarrow -1$;\quad $\beta^* \leftarrow -1$
    \FOR{$j = 1$ \textbf{to} $N$, $j\neq i$}
      \IF{$\text{fit}_j < \text{fit}_i$}
        \STATE $r \leftarrow \|\mathbf{x}_i - \mathbf{x}_j\|_2$;\quad
               $\beta \leftarrow \beta_0 e^{-\gamma_{\text{eff}}r^2}$
        \IF{$\beta > \beta^*$}
          \STATE $\beta^* \leftarrow \beta$;\; $j^* \leftarrow j$
        \ENDIF
      \ENDIF
    \ENDFOR
    \IF{$j^* \ge 0$}
      \STATE $\mathbf{x}_i^{\text{new}} \leftarrow
             \mathrm{clip}(\mathbf{x}_i + \beta^*(\mathbf{x}_{j^*}-\mathbf{x}_i)
             + \alpha(\mathbf{u}-0.5),\,\mathbf{lb},\,\mathbf{ub})$
    \ELSE
      \STATE $\mathbf{x}_i^{\text{new}} \leftarrow
             \mathrm{clip}(\mathbf{x}_i + 2\alpha(\mathbf{u}-0.5),\,\mathbf{lb},\,\mathbf{ub})$
    \ENDIF
  \ENDFOR
  \STATE $\mathbf{x} \leftarrow \mathbf{x}^{\text{new}}$;\quad
         $\text{fit}_i \leftarrow f(\mathbf{x}_i)$ for all $i$
  \STATE Update $\mathbf{x}^*$, $f^*$ if improved
\ENDFOR
\RETURN $\mathbf{x}^*$
\end{algorithmic}
\end{algorithm}

\textbf{Complexity Analysis}

\begin{itemize}
  \item \textbf{Time Complexity:} $O(G\cdot N^2\cdot D)$. The innermost double loop
        ($N\times N$ firefly comparisons) executes $G$ times. Each pairwise comparison
        computes a $D$-dimensional Euclidean distance in $O(D)$, giving $O(G\cdot N^2\cdot D)$
        for movement alone. Fitness re-evaluation costs $O(G\cdot N\cdot F)$.
  \item \textbf{Space Complexity:} $O(N\cdot D)$ for the position array and its copy
        (\texttt{new\_fireflies}).
\end{itemize}

\textbf{Advantages}

\begin{itemize}
  \item \textbf{Single-best-attractor rule reduces noise:} By moving toward only the
        \textit{most attractive} brighter neighbour (rather than a weighted sum of all of them),
        each update has a clear, consistent direction; this avoids the conflicting forces
        that arise when many bright fireflies pull in different directions simultaneously.
  \item \textbf{Automatic sub-swarm formation:} Because attractiveness decays with distance,
        nearby fireflies tend to cluster around shared local optima while distant sub-groups
        explore independently, providing natural multi-start behaviour on multimodal landscapes.
  \item \textbf{Scale-invariant absorption:} Normalising $\gamma$ by the squared domain
        diameter makes the algorithm's exploration radius consistent regardless of whether the
        search space spans $[-5,5]^D$ or $[-500,500]^D$.
  \item \textbf{Adaptive step decay:} The $0.97^t$ schedule requires no additional parameter
        and smoothly narrows the random perturbation from coarse exploration to fine
        exploitation without manual tuning.
\end{itemize}

\textbf{Disadvantages}

\begin{itemize}
  \item \textbf{Quadratic per-iteration cost:} The $O(N^2)$ pairwise scan makes FA
        substantially slower than PSO ($O(N)$) or the discrete ACO ($O(N\cdot n)$) at equal
        population sizes. With the default $N=50$, each iteration requires 2\,450 pairwise
        distance computations.
  \item \textbf{Poor scalability:} Doubling $N$ quadruples the movement cost. For large
        populations or high-dimensional problems, runtime becomes prohibitive.
  \item \textbf{Sensitivity to $\gamma$:} If $\gamma$ is too large, $\beta_{ij}\approx 0$ for
        all but the nearest firefly, turning the algorithm into a purely local search with very
        limited global exploration. If $\gamma$ is too small, all fireflies feel equal
        attraction regardless of distance, collapsing into PSO-like behaviour without the
        efficiency of PSO.
  \item \textbf{Collapse on narrow-valley functions:} On Rosenbrock, the attraction-based
        jumps ($\beta\cdot(\mathbf{x}_{j^*}-\mathbf{x}_i)$) tend to overshoot the narrow
        curved valley, pushing positions far from the optimum and causing objective values
        to diverge to extremely large numbers, as confirmed experimentally.
\end{itemize}
